\section*{Capítulo \ref{ch:g6:measurement-in-science} --- Medir en ciencia: unidades e instrumentos}

\begin{solution}{exo:g6:measurement-in-science:1}
$2$ \unit{m}; \ $8$ \unit{g}; \ $55$ \unit{min}; \ $39$
\unit{\celsius}.
\end{solution}

\begin{solution}{exo:g6:measurement-in-science:2}
\qty{320}{cm}; \ \qty{0.75}{kg}; \ \qty{240}{s}; \ \qty{5.6}{km}.
\end{solution}

\begin{solution}{exo:g6:measurement-in-science:3}
Sus peldaños van de $60$ en $60$ y no de $10$ en $10$, así que la
escalera decimal miente:
$\qty{3.5}{min} = 3 \times 60 + 30 = \qty{210}{s}$, mientras que la
escalera habría dado, mal, \qty{350}{s}.
\end{solution}

\begin{solution}{exo:g6:measurement-in-science:4}
Cada paso: $200 \div 5 = \qty{40}{mL}$. Dos pasos por encima del
$300$: $300 + 80 = \qty{380}{mL}$.
\end{solution}

\begin{solution}{exo:g6:measurement-in-science:5}
Levadura: la báscula de cocina (pasos de \omterm{def:g3:measuring-mass:units}{gramo} --- la del baño no
notaría una pizca). Latidos: un cronómetro (pasos de segundo y más
finos). Ancho de la página: una regla (pasos de milímetro). Patio del
colegio: una cinta métrica larga o una rueda de \omterm{def:g6:measurement-in-science:measuring}{medir} (alcance de
varios \omterm{def:g2:measuring-length:units}{metros}).
\end{solution}

\begin{solution}{exo:g6:measurement-in-science:6}
El paso de juzgar: un lápiz se estima en menos de \qty{20}{cm}, así
que \qty{17.4}{m} suspende el examen de sensatez al instante. Leyó
\omterm{def:g2:measuring-length:units}{centímetros} y escribió \omterm{def:g2:measuring-length:units}{metros}: el valor verdadero es
\qty{17.4}{cm}.
\end{solution}

\begin{solution}{exo:g6:measurement-in-science:7}
Honrada: ``entre $18$ y \qty{19}{\celsius}, más cerca del $19$'' (o
``unos \qty{18.8}{\celsius}''). Deshonesta:
``\qty{18.7643}{\celsius}'' --- una precisión que las marcas no pueden
prometer.
\end{solution}

\begin{solution}{exo:g6:measurement-in-science:8}
No tiene la culpa nadie: los dedos arrancan y paran con un pelo de
diferencia, y una dispersión pequeña es la respiración normal de la
medida. Manera sabia de comunicarlo: las lecturas se agrupan en torno
a \qty{6.2}{s} --- ``unos \qty{6.2}{s}''.
\end{solution}

\begin{solution}{exo:g6:measurement-in-science:9}
Un cuarto de \qty{450}{g}: unos \qty{112}{g} (\qty{113}{g} también
vale). El recado enseña por qué importan las \omterm{def:g6:measurement-in-science:measuring}{unidades} compartidas:
las recetas escritas con las \omterm{def:g6:measurement-in-science:measuring}{unidades} de un país hay que
traducirlas antes de que los instrumentos de otro país puedan
obedecerlas.
\end{solution}

\begin{solution}{exo:g6:measurement-in-science:10}
El sello va de la marca del $1$ a la del $3.4$:
$3.4 - 1 = \qty{2.4}{cm}$. La costumbre: no fiarse nunca del borde de
un instrumento --- se empieza en una marca elegida y se resta, y así
las reglas gastadas pierden su poder.
\end{solution}

\begin{solution}{exo:g6:measurement-in-science:11}
Las dos lecturas se separan \qty{0.10}{m}, así que la última cifra de
cualquiera de las dos no es de fiar; se agrupan en torno a $18.3$.
Escribir ``\qty{18.3}{m}'' conserva todas las cifras que la cinta
puede prometer con honradez y ninguna más.
\end{solution}

\begin{solution}{exo:g6:measurement-in-science:12}
Desastres: los ``pies'' de dos pueblos no coinciden nunca, así que las
vigas cortadas en un pueblo no encajan jamás en las casas de otro; y,
cuando hereda el hijo, la \omterm{def:g6:measurement-in-science:measuring}{unidad} misma cambia de longitud --- todas
las medidas del reino caducan con el rey. La barra de metal arregló
las dos cosas: una longitud que no cambia, copiada para todos los
pueblos y que sobrevive a todos los reyes. Su propio defecto: un
objeto único se puede rayar, robar o alterar poco a poco, y el mundo
entero tenía que peregrinar hasta su cámara acorazada para comprobar
las copias --- hasta que las \omterm{def:g6:measurement-in-science:measuring}{unidades} se reconstruyeron sobre las
constantes de la naturaleza, que no se desgastan y están de servicio
en todas partes a la vez.
\end{solution}

\begin{solution}{pb:g6:measurement-in-science:1}
\textbf{1.} Cada pesada suma la carga a lo que la aguja ya marcaba;
un arranque distinto de cero envenena todas las lecturas del día ---
la comprobación del arranque justo tiene que ir primero.
\textbf{2.} $500 + 200 = \qty{700}{g}$.
\textbf{3.} Marca \qty{20}{g} de menos. Echando harina hasta que la
aguja dice \qty{700}{g}, el panadero ha echado en realidad
\qty{720}{g}: la báscula hace que el panadero entregue más mercancía
de la que paga el cliente --- favorece a los clientes, y el panadero
pierde en cada venta.
\textbf{4.} $500 + 200 + 50 = \qty{750}{g}$.
\textbf{5.} No: un sello de \qty{2}{g} desaparece por debajo de los
pasos de \qty{5}{g} de la báscula --- está fuera de su territorio.
\textbf{6.} $200 \div 4 = \qty{50}{mL}$ por paso.
\textbf{7.} El agua está en $600 + 50 = \qty{650}{mL}$; ``unos
\qty{700}{mL}'' redondea y se come un paso entero --- demasiado vago
para un instrumento tan bueno.
\textbf{8.} El ojo a la altura de la marca: leer desde abajo desplaza
la línea del agua aparente --- se incumplió la regla de mirar justo
enfrente.
\textbf{9.} Se equivoca en \qty{10}{mL} --- mejor (más pequeño) que su
propio paso de \qty{50}{mL}: aceptable para una jarra de colegio.
\textbf{10.} Adelanta, $2$ minutos.
\textbf{11.} $2 \times 7 = 14$ minutos de adelanto al cabo de una
semana.
\textbf{12.} Primer reloj: atrasarlo $2$ min; segundo: adelantarlo
$2$ min; tercero: nada --- está bien. (Promediar los tres no
``arreglaría'' nada y falsearía el único reloj honrado.)
\textbf{13.} Por ejemplo: \emph{Cero}: comprueba que el instrumento
vacío marca cero antes de fiarte de nada de lo que diga.
\emph{Lectura}: aprende cuánto vale el paso y lee justo enfrente, y
entre marcas con modestia honrada. \emph{\omterm{def:g6:measurement-in-science:measuring}{Unidades}}: un número sin
su \omterm{def:g6:measurement-in-science:measuring}{unidad} no es una medida --- escribe la \omterm{def:g6:measurement-in-science:measuring}{unidad}, siempre.
\end{solution}
